\documentclass[11pt]{article}
\usepackage[margin=1in]{geometry}
\usepackage{amsmath,amssymb,amsthm,mathtools}
\usepackage[hidelinks]{hyperref}
\newtheorem{theorem}{Theorem}
\newtheorem{lemma}[theorem]{Lemma}
\newtheorem{corollary}[theorem]{Corollary}
\title{All ordinary one-variable constant-plus-power Strong Factorial families}
\author{Internal uniform contour proof and exact finite certificates; no priority claim}
\date{17 November 2026}
\begin{document}
\maketitle
Put
\[
 F_{n,m}(t)=\sum_{k=0}^n\binom nk(mk)!t^k,\qquad
 P_{n,m}(z)=z^nF_{n,m}(1/z),\qquad A_m=(m/e)^m.
\]
All Gamma shapes are one and $n,m$ are positive integers. The goal of
resolving general Factorial and/or unrestricted Strong Factorial remains
open: this note concerns precisely $A+BY^m$ in one positive exponential
variable. It does not assert a reduction from arbitrary polynomial supports.

\begin{theorem}\label{thm:all}
For every $n,m\ge1$, $F_{n,m}$ and $P_{n,m}$ have simple complex roots.
Consequently $A+BY^m$ satisfies Strong Factorial at every starting index,
for every positive integer $m$ and every nonzero such complex polynomial.
\end{theorem}
The analytic cutoff is $n=10000$, uniformly over all exponents after using
the elementary large-exponent criterion. The finite part reuses the
registered prime-power and exponent-six results and verifies all other
necessary pairs below this cutoff exactly. No density theorem is used.

\section{The uniform analytic theorem}
\begin{theorem}\label{thm:analytic}
Let $n\ge10000$ and $6\le m\le\sqrt n$. Every root satisfies
\[
             \frac12<\left|\frac{z}{n^m A_m}\right|<\frac{21}{10},
                                                               \tag{1}
\]
and every root is simple. The estimates below are uniform in both integers.
\end{theorem}
For $v=z/(n^m A_m)$, set
\[
 I_{n,m}(v)=\int_0^\infty(A_mv+s^m)^ne^{-ns}\,ds
           =K_{n,m}(v)+J_{n,m}(v),                    \tag{2}
\]
where the two intervals are $[0,m/2]$ and $[m/2,\infty)$.
Then $P_{n,m}(n^mA_mv)=n^{mn+1}I_{n,m}(v)$.
Both pieces in (2) are polynomials in $v$; the decomposition is exact.
We abbreviate them as $K,J$, and set $B=11/5$, $R=21/10$, $\delta=1/25$.

\subsection{The outer coefficient bound}
For $\theta=k/n$, elementary entropy and factorial integral estimates give
\[
 \frac{[\binom nk(mk)!]^{1/k}}{n^mA_m}
 \le q_m(\theta)(emk)^{1/k},\quad
 q_m(\theta)=\theta^{m-1}(1-\theta)^{-(1-\theta)/\theta}.
                                                               \tag{3}
\]
The endpoints are continuous limits. Since $m\ge6$, $q_m\le q_6<300/299$.
The last bound is proved in the
\href{EFFECTIVE_EXPONENT_SIX_SIMPLICITY.tex}{exponent-six proof}:
$\log q_6$ has derivative numerator $6\theta+\log(1-\theta)$;
its maximum is at $\theta=1-u$, $0<u<1/300$, where
$\log q_6=5\log(1-u)+6u<u$.

If $k\ge n/10$, $m\le\sqrt n$ yields
\[
 \log (emk)^{1/k}\le\frac{10(1+\tfrac32\log n)}n
 <\frac2{125}\quad(n\ge10000).
\]
The expression decreases in $n$, and $\log10000<10$.
Thus $(emk)^{1/k}<125/123$, and
$(300/299)(125/123)<21/20$.
If $k<n/10$, use instead $\binom nk^{1/k}\le en/k$ and
$(emk)^{1/k}\le em$ to bound the left side of (3) by
$e^2m\theta^{m-1}<9m/10^{m-1}\le54/10^5<21/20$.
For the elementary inequality just used, $k\le2^{k-1}\le(em)^{k-1}$.
Every coefficient therefore has $k$th root less than
$(21/20)n^mA_m$. A finite geometric-series bound excludes every
$|z|\ge Rn^mA_m$.

\subsection{A single large saddle on a full disk}
For every $|v|\le B$, the equation
\[
                 A_mv=s^{m-1}(m-s)                    \tag{4}
\]
has a unique solution $s_0=\alpha+i\beta$ with $|s_0-m|<\delta$.
On the circle, the ratio of the lower bound to $A_m$ is
\[
 \delta\frac{e^m}m(1-\delta/m)^{m-1}
 \ge\delta(1-\delta)\frac{e^m}m>B.
\]
Bernoulli's inequality supplies the middle bound. The last expression
increases for $m\ge6$, and at six the rational replacement $e>27/10$
already proves the strict inequality. Rouch\'e therefore applies.
The saddle is simple and depends analytically on $v$ on a neighborhood
of the closed disk. Near it, use
\[
 \Phi(s,v)=m\log s+\log(1+A_mv/s^m)-s,\quad
 \Phi_0=\Phi(s_0,v),\quad a=1-(m-1)/s_0.
\]
Then $\Phi_s(s_0,v)=0$, $\Phi_{ss}(s_0,v)=-a$, and
\[
 \operatorname{Re}a>\frac9{10m},\qquad |a|<\frac{11}{10m}. \tag{5}
\]
These follow from $(1-\delta)/(m-\delta)$ and
$(1+\delta)/(m-\delta)$, respectively.
Put $G_m(x)=m\log x-x$ and $S=\operatorname{Re}\Phi_0$.
Taylor's inequality for $\log(1+u)$ gives
\[
 S>G_m(m)-1/100=\log A_m-1/100,                        \tag{6}
\]
since the loss is at most
$\delta/m+\delta^2/[2m(1-\delta/6)]$, whose maximum at $m=6$ is
less than $1/100$.

\subsection{The complete uniform contour estimate}
\begin{lemma}\label{lem:J}
For $n\ge10000,m\ge6$, on $|v|\le B$,
\[
 J(v)=U(v)(1+\varepsilon(v)),\quad
 U=e^{n\Phi_0}\sqrt{2\pi/(na)},\quad
 |\varepsilon|<\frac{22}{\sqrt{mn}}<\frac1{10}.         \tag{7}
\]
On $|v|\le R$, this gives
\[
 \left|\frac{J'}J-n\frac{m-s_0}{mv}\right|<3.          \tag{8}
\]
No condition $m\le\sqrt n$ is needed for this lemma.
\end{lemma}
\begin{proof}
Deform $[m/2,\infty)$ to its vertical translate by $i\beta$, keeping the
vertical segment at $m/2$. The integrand is entire and the closing side
at positive infinity tends to zero. On the local horizontal interval
$s=s_0+t$, $|t|\le r=m/8$, we have $|s|>6m/7$.
Put $q_t=B(35/81)^6$. Differentiating the absolutely convergent
$\log(1+A_mv/s^m)$ series shows that
\[
 m^2|\Phi_{sss}|
 \le(7/6)^3\left[2+
 \frac{36q_t(1+q_t)}{(1-q_t)^3}
 +\frac{18q_t}{(1-q_t)^2}+\frac{2q_t}{1-q_t}\right]<8. \tag{9}
\]
Indeed $|A_mv/s^m|\le B(7/(6e))^m\le B(35/81)^m$.
The quantities $m^2(35/81)^m$, $m(35/81)^m$ and $(35/81)^m$
decrease for $m\ge6$, so their value at six bounds the derivative series.

The Taylor remainder is at most $4|t|^3/(3m^2)$.
Using (5), the difference between the local integral and its Gaussian is
at most
\[
 \frac{4ne^{nS}}{3m^2}\int_{\mathbb R}|t|^3
 e^{-17nt^2/(60m)}dt
 =\frac{4800e^{nS}}{289n}<\frac{17e^{nS}}n.
\]
Completing the Gaussian costs at most
$[160/(9n)]e^{nS-9mn/1280}<18e^{nS-9mn/1280}/n$.
Its complete value $U$ has modulus $>2\sqrt{m/n}\,e^{nS}$ by (5).

We bound every remaining contour piece. On the horizontal line
$s=x+i\beta$, $2m/3\le x\le4m/3$, put $q_c=B(5/9)^6$.
The real second derivative satisfies
\[
 \frac{d^2}{dx^2}\operatorname{Re}\Phi(x+i\beta,v)
 \le-\frac m{x^2}\left[
 \frac{1-10^{-4}}{(1+10^{-4})^2}
 -\frac{6q_c}{(1-q_c)^2}-\frac{q_c}{1-q_c}\right]
 <-\frac m{3x^2}\le-\frac3{16m}.                      \tag{10}
\]
Here $|\beta|/x\le1/100$,
$|A_mv/s^m|\le B(3/(2e))^m\le B(5/9)^m$, and
$m(5/9)^m$ decreases for $m\ge6$.
The first derivative vanishes at $x=\alpha$. Thus the height is at most
$S-3(x-\alpha)^2/(32m)$. Outside $|x-\alpha|\le m/8$, its loss is
$>m/1000$. The length is less than $m$, giving
$m e^{nS-mn/1000}$.
For $x\ge4m/3$, put $q_f=B(5/18)^6$; the derivative is bounded by
\[
 \frac{m}{x(1-|A_mv/s^m|)}-1
 \le\frac{3/4}{1-q_f}-1<-1/5.
\]
At $4m/3$ the height loss in (10) is at least
$(3m/32)(49/150)^2>m/100$. The infinite tail is therefore at most
$5e^{nS-mn/100}/n$.

On the interval $m/2\le x\le2m/3$, the upper envelope
$[BA_m+(x+\delta)^m]e^{-x}$ increases: its derivative has the sign of
$u^{m-1}(m-u)-BA_m>0$ for $u=x+\delta<m-1$.
The latter map is increasing there and its value at $m/2$ already
exceeds $BA_m$, since $(e/2)^m>B$ for $m\ge6$.
Relative to $A_m$, the envelope at $2m/3$ is at most
\[
 B e^{-2m/3}+r_\ell^m<e^{-m/50},\qquad
 r_\ell=(101/150)(141/100).
\]
This uses $e^{1/3}<141/100$ and $\delta/m\le1/150$.
For an exact check valid for all $m\ge6$, multiply by $e^{m/50}$;
both resulting bases are less than one, and their value at six is bounded by
$B(10/27)^3+r_\ell^6/(22/25)<1$.
By (6), the height is below $S-m/100$; this piece has length $m/6$.
For the vertical segment, the corresponding upper bound relative to
$A_m$ is
\[
 B e^{-m/2}+r_v^m<e^{-m/50},\qquad r_v=38/45,
\]
using $e^{1/2}<5/3$. The same decreasing-bases argument at six gives
$B(10/27)^2+r_v^6/(22/25)<1$. Its length is at most $\delta$, and its
height is below $S-m/100$.
All logarithmic charts on the horizontal part are valid: for $x\ge m/2$,
$|A_mv/s^m|\le B(2/e)^m<1$ already at $m=6$ and then decreases.

Combining all pieces, the error after removing $U$ is at most $e^{nS}$
times
\[
 \frac{40}{n}+m e^{-mn/1000}+(m+\delta)e^{-mn/100}
 <\frac{44}{n}\quad(n\ge10000,m\ge6).                \tag{11}
\]
Use $e^u\ge u^3/6$ on the first exponential and $e^u\ge u^2/2$ on
the second. After multiplication by $n$, the upper constant is
$40+6\cdot1000^3/(36\cdot10000^2)
 +2\cdot100^2(6+\delta)/(36\cdot10000)<44$.
This proves (7).

The relative error is analytic on a larger disk, since $J$ is a polynomial
and $U$ is analytic and nonzero. Radius-$1/10$ Cauchy circles give
$|\varepsilon'|\le220/\sqrt{mn}<1$ on $|v|\le R$.
Also
\[
 \frac{U'}U=n\frac{m-s_0}{mv}-\frac{a'}{2a},\qquad
 s_0'=\frac{A_m}{m s_0^{m-2}(m-1-s_0)}.
\]
Bernoulli's inequality gives
$|s_0'|\le me^{-m}/(1-\delta)^2<1$, uniformly for $m\ge6$.
Equation (5) then gives $|a'/(2a)|<1$.
The logarithmic derivative correction is below $1+10/9<3$, proving (8).
The removable value at $v=0$ is understood in the phase formula.
\end{proof}

\subsection{The endpoint, root annulus and simplicity}
\begin{lemma}\label{lem:K}
For $m\ge6,n\ge10000$ and $|v|\ge2/5$,
\[
 K(v)=\frac{(A_mv)^n}{n}(1+D(v)),\qquad
 |D|<1/1000,\quad |D'|<1/1000,\quad
 \left|K'/K-n/v\right|<1.                            \tag{12}
\]
\end{lemma}
\begin{proof}
For $0\le x\le m/2$,
\[
 \log\left(1+\frac{x^m}{(2/5)A_m}\right)\le19x/20.     \tag{13}
\]
At the right endpoint this is equivalent to
$1+(5/2)(e/2)^m<e^{19m/40}$.
After multiplication by $e^{-19m/40}$, both terms decrease with $m$.
At $m=6$ the inequality is certified by
\[
 1+(5/2)(49/36)^6<\sum_{j=0}^{12}\frac{(57/20)^j}{j!}<e^{57/20},
\]
using $e<49/18$.
The second base $e^{21/40}/2$ is less than one, since
$e^{21/40}\le(5/3)(40/39)<2$.
To extend (13) to the full interval, the derivative of the logarithm
divided by $x$ has the sign of
$h(u)=mu/(1+u)-\log(1+u)$.
It increases on $0\le u\le m-1$. For larger $u$ in the relevant range,
$mu/(1+u)\ge m-1>19m/40\ge\log(1+u)$.
Thus that quotient increases and its endpoint bounds it.

As in the exponent-six proof, subtracting the full exponential integral
and using $|(1+z)^n-1|\le n|z|(1+|z|)^{n-1}$ gives
\[
 |D|\le\frac{m!20^{m+1}}{(2/5)A_m n^{m-1}}+e^{-mn/2},\qquad
 |D'|\le\frac{m!20^{m+1}}{(2/5)^2A_m n^{m-1}}.
\]
The elementary bound $m!\le em(m/e)^m$ bounds these expressions by
$150mn(20/n)^m+e^{-mn/2}$ and $375mn(20/n)^m$.
Both are maximized at $m=6,n=10000$ and are less than $1/1000$.
The logarithmic derivative claim follows.
\end{proof}

For $|v|\le1/2$, the absolute bound is
$|K|\le(m/2)(A_m/2)^n$. Indeed the maximum of
$[A_m/2+x^m]e^{-x}$ on $[0,m/2]$ occurs at an endpoint.
The right endpoint is below $A_m/2$, since
\[
 e^{-m/2}+2(e^{1/2}/2)^m
 <(5/8)^m+2(5/6)^m\le(5/8)^6+2(5/6)^6<1.
\]
By (6)--(7), $|J|>(3/2)\sqrt{m/n}\,(A_me^{-1/100})^n$.
Therefore
\[
 \frac{|K|}{|J|}<\frac{\sqrt{mn}}3(50/99)^n<1,
\]
when $m\le\sqrt n$ and $n\ge10000$; use $\sqrt{mn}\le n$ and the
exact comparison $10000(50/99)^{10000}<3$, decreasing thereafter.
This excludes the inner disk. Combined with (3), it proves (1).

On that entire annulus, $J$ and $K$ are nonzero. At a zero of $I=J+K$,
\[
 I'=J(J'/J-K'/K),\qquad
 \left|\frac{J'}J-\frac{K'}K\right|
 >\frac{n|s_0|}{m|v|}-4>\frac n3-4>0.                \tag{14}
\]
Here $|s_0|/m>149/150$ and $|v|<21/10$.
Thus every possible zero is simple, proving Theorem~\ref{thm:analytic}.
The exact interval split and the full-disk contour estimate replace any
unproved assumption about selecting formal saddle contributions.

\section{Covering every exponent and the complete finite prefix}
The previously proved
\href{ALL_EXPONENTS_AND_LARGE_EXPONENT_CUTOFF.tex}{coefficient criterion}
states $2^m\ge4n^2\Rightarrow F_{n,m}$ squarefree.
If $n\ge10000$ and it does not apply, then
\[
 m<2\log_2(2n)<\sqrt n.
\]
For the last inequality put $u=\sqrt n\ge100$:
$u-2-4\log_2u$ is positive at one hundred and increases, using
$\log2>1/2$ and $\log_2100<7$.
Exponents $1\le m\le5$ are already covered by the retained
\href{PRIME_POWER_TWO_TERM_SFC.tex}{prime-power theorem}, including one.
Theorem~\ref{thm:analytic} handles all other exponents at these indices.
Thus the analytic cutoff $10000$ covers every positive exponent.

For $1\le n<10000$, $m\ge30$ satisfies the coefficient criterion;
$n=1$ is linear at every exponent. Prime powers and $m=1$ are already
settled, and the full
\href{EFFECTIVE_EXPONENT_SIX_SIMPLICITY.tex}{exponent-six result} is reused.
The only further necessary exponents are
\[
                  10,12,14,15,18,20,21,22,24,26,28.
\]
There are exactly $92935$ pairs in this range with $2^m<4n^2$.
The \href{all_exponent_prefix_certificates.json}{finite certificate file}
covers every one. For $92900$ pairs, take the preceding prime $q$ and
$d=n-q$; the file verifies $q>m\max(1,d)$ and gives a modular B\'ezout
identity for $F_{d,m}$ and its derivative with right side of degree at most
one. The registered
\href{QUANTITATIVE_SQUAREFREENESS_ALL_EXPONENTS.tex}{shifted Newton lemma}
forces a rational square divisor to reduce into this kernel preserving its
degree. The kernel degree is at most $35$.
For the remaining $35$ pairs, a prime $p>mn$ gives the same certificate
for the full-degree reduction of $F_{n,m}$; their largest index is $126$.
Every selected prime and every B\'ezout identity is checked exactly.

Finally, any repeated irreducible rational factor has even degree.
For positive exponential $Y$, $F_{n,m}(t)=\mathbb E(1+tY^m)^n$.
Even $n$ gives no real root; odd $n$ has strictly positive derivative on
the real axis, so all real roots are simple. An odd-degree repeated
rational factor would have a real root. Hence a nonconstant rational
square divisor has degree at least two, contradicting the degree-at-most-one
B\'ezout identity after degree-preserving reduction. Every complex repeated
root supplies a repeated rational factor in characteristic zero.
This proves the complete finite prefix and Theorem~\ref{thm:all}.
The corollary follows from $P'_{n+1,m}=(n+1)P_{n,m}$ and
$\mathcal L((A+BY^m)^n)=B^nP_{n,m}(A/B)$, with one-term cases immediate.

\section{Scope and verification boundary}
For this ordinary constant-plus-power sequence the exceptional set is now
empty, so $E_m(X)=D_{\rm all}(X)=0$. The earlier density estimates and
source-qualified analytic-number-theory routes remain retained evidence,
but are unnecessary for this direct pointwise result.
This is not unrestricted Strong Factorial: both-nonconstant monomial
supports, general polynomial supports and other Gamma shapes are not
asserted here. In particular it is not a proof of general Factorial.
The checker verifies rational constants, algebraic identities and every
finite-prefix B\'ezout certificate. The universal evidence is the full
analytic proof. No numerical zero scan, priority claim, independent replay,
formal verification or external review is claimed.
\end{document}
